% GENERATED FILE. Edit canonical lessons, then run npm run generate:books.
% canonical-source-hash: fnv1a64:95ca10f9c303a80b
% canonical-lessons: MR-C386-pratibandh-karne, MR-C386-rokhne, MR-C386-adavne, MR-C386-dakhal-ghene, MR-C386-palvat

\chapter{Stopping Trouble}
\label{ch:mr-stopping-trouble}
\hlchaptermodality{386}

\begin{chapteropening}
\textbf{By the end of this chapter:} \emph{I can talk about preventing, blocking and noticing trouble.}

Complete the last lesson of chapter 386: I can talk about preventing, blocking and noticing trouble.
\end{chapteropening}

\section[\texorpdfstring{\mr{प्रतिबंध} \mr{करणे} (pratibandh karṇe)}{प्रतिबंध करणे (pratibandh karṇe)}]{\mr{प्रतिबंध} \mr{करणे} (pratibandh karṇe) — to prevent}

\label{lesson:MR-C386-pratibandh-karne}



\begin{quote}
Before the new one: say the Marathi for to be averted, then the Marathi for to confront.
\end{quote}

\subsection*{You'll want to know: \mr{प्रतिबंध} \mr{करणे}}
\textbf{\mr{प्रतिबंध} \mr{करणे}} — \emph{pratibandh karṇe} — \textquotedblleft{}to prevent\textquotedblright{}.

\begin{quote}
\mr{स्वच्छतेमुळे} \mr{आजारांना} \mr{प्रतिबंध} \mr{करता} \mr{येतो}.
\end{quote}

\begin{quote}
\emph{svacchatemuḷe ājārānnā pratibandh kartā yeto.}
\end{quote}

\begin{quote}
Cleanliness can prevent illness.
\end{quote}

Past, with -\mr{ने} (he …): \mr{त्याने} \mr{प्रतिबंध} \mr{केला}, \emph{tyāne pratibandh kelā}.

What is prevented takes -\mr{ला}; \mr{प्रतिबंधक}, preventive.

\begin{cousinweb}
\textbf{\mr{प्रतिबंध} \mr{करणे}} (\emph{pratibandh karṇe}) joins the tatsama \mr{प्रतिबंध}, Sanskrit \emph{pratibandha}, \textquotedblleft{}connexion, hindrance\textquotedblright{} (CDIAL 8580), to \mr{करणे}, \textquotedblleft{}to do\textquotedblright{}.
\end{cousinweb}

\begin{grammarlens}[title={What you've built}]
One more word for saying what happens if.
\end{grammarlens}

\subsection*{Your turn}
\begin{itemize}
\raggedright
  \item \emph{Say it:} \emph{pratibandh karṇe}
  \item \emph{Say it:} \emph{pratibandh karṇe}, once more
  \item \emph{Say it:} \emph{svacchatemuḷe ājārānnā pratibandh kartā yeto.} — Cleanliness can prevent illness.
  \item \emph{Recall:} say the Marathi for to be averted, then the Marathi for to confront, then say \emph{pratibandh karṇe} again
\end{itemize}

\begin{tcolorbox}[breakable,colback=teal!4,colframe=teal!35!black,title={Before you move on}]
What does \mr{प्रतिबंध} \mr{करणे} mean? (\textquotedblleft{}To prevent\textquotedblright{}.) Say it once more.
\end{tcolorbox}
\section[\texorpdfstring{\mr{रोखणे} (rokhṇe)}{रोखणे (rokhṇe)}]{\mr{रोखणे} (rokhṇe) — to restrain; to stop}

\label{lesson:MR-C386-rokhne}



\begin{quote}
Before the new one: say the Marathi for to confront, then the Marathi for to prevent.
\end{quote}

\subsection*{You'll want to know: \mr{रोखणे}}
\textbf{\mr{रोखणे}} — \emph{rokhṇe} — \textquotedblleft{}to restrain; to stop\textquotedblright{}.

\begin{quote}
\mr{पोलिसांनी} \mr{गर्दीला} \mr{दरवाज्याजवळ} \mr{रोखलं}.
\end{quote}

\begin{quote}
\emph{polisānnī gardīlā darvājyājavaḷ rokhlaṁ.}
\end{quote}

\begin{quote}
The police held the crowd back near the gate.
\end{quote}

Past, with -\mr{ने} (he …): \mr{त्याने} \mr{रोखलं}, \emph{tyāne rokhlaṁ}.

What is stopped takes -\mr{ला}: \mr{गर्दीला} \mr{रोखणे}.

\begin{cousinweb}
\textbf{\mr{रोखणे}} (\emph{rokhṇe}) is inherited: CDIAL lists Marathi \emph{rokhṇẽ}, \textquotedblleft{}to hinder\textquotedblright{}, under \emph{*rōkk-}, \textquotedblleft{}stop\textquotedblright{} (CDIAL 10827).
\end{cousinweb}

\begin{grammarlens}[title={What you've built}]
One more word for saying what happens if.
\end{grammarlens}

\subsection*{Your turn}
\begin{itemize}
\raggedright
  \item \emph{Say it:} \emph{rokhṇe}
  \item \emph{Say it:} \emph{rokhṇe}, once more
  \item \emph{Say it:} \emph{polisānnī gardīlā darvājyājavaḷ rokhlaṁ.} — The police held the crowd back near the gate.
  \item \emph{Recall:} say the Marathi for to confront, then the Marathi for to prevent, then say \emph{rokhṇe} again
\end{itemize}

\begin{tcolorbox}[breakable,colback=teal!4,colframe=teal!35!black,title={Before you move on}]
What does \mr{रोखणे} mean? (\textquotedblleft{}To restrain; to stop\textquotedblright{}.) Say it once more.
\end{tcolorbox}
\section[\texorpdfstring{\mr{अडवणे} (aḍavṇe)}{अडवणे (aḍavṇe)}]{\mr{अडवणे} (aḍavṇe) — to block; to stop (someone)}

\label{lesson:MR-C386-adavne}



\begin{quote}
Before the new one: say the Marathi for to prevent, then the Marathi for to restrain.
\end{quote}

\subsection*{You'll want to know: \mr{अडवणे}}
\textbf{\mr{अडवणे}} — \emph{aḍavṇe} — \textquotedblleft{}to block; to stop (someone)\textquotedblright{}.

\begin{quote}
\mr{रस्त्यात} \mr{एका} \mr{माणसाने} \mr{मला} \mr{अडवलं}.
\end{quote}

\begin{quote}
\emph{rastyāt ekā māṇsāne malā aḍavlaṁ.}
\end{quote}

\begin{quote}
A man stopped me on the road.
\end{quote}

Past, with -\mr{ने} (he …): \mr{त्याने} \mr{अडवलं}, \emph{tyāne aḍavlaṁ}.

The one stopped takes -\mr{ला}; \mr{अडणे}, to get stuck.

\begin{cousinweb}
\textbf{\mr{अडवणे}} (\emph{aḍavṇe}) is built on \mr{अडणे}, \textquotedblleft{}to be stopped\textquotedblright{}, as its causative; CDIAL lists it as Marathi \emph{aḍṇẽ} under \emph{*aḍḍ-}, \textquotedblleft{}obstruct, stop\textquotedblright{} (CDIAL 188).
\end{cousinweb}

\begin{grammarlens}[title={What you've built}]
One more word for saying what happens if.
\end{grammarlens}

\subsection*{Your turn}
\begin{itemize}
\raggedright
  \item \emph{Say it:} \emph{aḍavṇe}
  \item \emph{Say it:} \emph{aḍavṇe}, once more
  \item \emph{Say it:} \emph{rastyāt ekā māṇsāne malā aḍavlaṁ.} — A man stopped me on the road.
  \item \emph{Recall:} say the Marathi for to prevent, then the Marathi for to restrain, then say \emph{aḍavṇe} again
\end{itemize}

\begin{tcolorbox}[breakable,colback=teal!4,colframe=teal!35!black,title={Before you move on}]
What does \mr{अडवणे} mean? (\textquotedblleft{}To block; to stop (someone)\textquotedblright{}.) Say it once more.
\end{tcolorbox}
\section[\texorpdfstring{\mr{दखल} \mr{घेणे} (dakhal gheṇe)}{दखल घेणे (dakhal gheṇe)}]{\mr{दखल} \mr{घेणे} (dakhal gheṇe) — to take notice (of)}

\label{lesson:MR-C386-dakhal-ghene}



\begin{quote}
Before the new one: say the Marathi for to restrain, then the Marathi for to block.
\end{quote}

\subsection*{You'll want to know: \mr{दखल} \mr{घेणे}}
\textbf{\mr{दखल} \mr{घेणे}} — \emph{dakhal gheṇe} — \textquotedblleft{}to take notice (of)\textquotedblright{}.

\begin{quote}
\mr{शेवटी} \mr{प्रशासनाने} \mr{आमच्या} \mr{तक्रारीची} \mr{दखल} \mr{घेतली}.
\end{quote}

\begin{quote}
\emph{śevṭī praśāsnāne āmacyā takrārīcī dakhal ghetlī.}
\end{quote}

\begin{quote}
In the end the administration took notice of our complaint.
\end{quote}

Past, with -\mr{ने} (he …): \mr{त्याने} \mr{दखल} \mr{घेतली}, \emph{tyāne dakhal ghetlī}.

What is noticed takes -\mr{ची}: \mr{तक्रारीची} \mr{दखल} \mr{घेणे}.

\begin{cousinweb}
\textbf{\mr{दखल} \mr{घेणे}} (\emph{dakhal gheṇe}) joins \mr{दखल}, \textquotedblleft{}notice, interference\textquotedblright{}, a loan from Persian usually traced to Arabic \emph{dakhl}, \textquotedblleft{}entering\textquotedblright{}, to \mr{घेणे}, \textquotedblleft{}to take\textquotedblright{}.
\end{cousinweb}

\begin{grammarlens}[title={What you've built}]
One more word for saying what happens if.
\end{grammarlens}

\subsection*{Your turn}
\begin{itemize}
\raggedright
  \item \emph{Say it:} \emph{dakhal gheṇe}
  \item \emph{Say it:} \emph{dakhal gheṇe}, once more
  \item \emph{Say it:} \emph{śevṭī praśāsnāne āmacyā takrārīcī dakhal ghetlī.} — In the end the administration took notice of our complaint.
  \item \emph{Recall:} say the Marathi for to restrain, then the Marathi for to block, then say \emph{dakhal gheṇe} again
\end{itemize}

\begin{tcolorbox}[breakable,colback=teal!4,colframe=teal!35!black,title={Before you move on}]
What does \mr{दखल} \mr{घेणे} mean? (\textquotedblleft{}To take notice (of)\textquotedblright{}.) Say it once more.
\end{tcolorbox}
\section[\texorpdfstring{\mr{पळवाट} (paḷvāṭ)}{पळवाट (paḷvāṭ)}]{\mr{पळवाट} (paḷvāṭ) — a loophole; an escape route}

\label{lesson:MR-C386-palvat}



\begin{quote}
Before the new one: say the Marathi for to block, then the Marathi for to take notice.
\end{quote}

\subsection*{You'll want to know: \mr{पळवाट}}
\textbf{\mr{पळवाट}} — \emph{paḷvāṭ} — \textquotedblleft{}a loophole; an escape route\textquotedblright{}.

\begin{quote}
\mr{या} \mr{नियमात} \mr{एक} \mr{पळवाट} \mr{आहे}.
\end{quote}

\begin{quote}
\emph{yā niymāt ek paḷvāṭ āhe.}
\end{quote}

\begin{quote}
There is a loophole in this rule.
\end{quote}

Feminine.

\mr{पळ} + \mr{वाट}, flight + path; \mr{पळवाट} \mr{काढणे}, to find a way out.

\begin{cousinweb}
\textbf{\mr{पळवाट}} (\emph{paḷvāṭ}) is a Marathi compound built from the stem of \mr{पळणे}, \textquotedblleft{}to run away\textquotedblright{}, inherited from Sanskrit \emph{palāyate}, \textquotedblleft{}flees\textquotedblright{} (CDIAL 7955), and \mr{वाट}, \textquotedblleft{}path\textquotedblright{}, inherited from Sanskrit \emph{vartman}, \textquotedblleft{}track of a wheel, path\textquotedblright{} (CDIAL 11366): a path to run by.
\end{cousinweb}

\begin{grammarlens}[title={What you've built}]
One more word for saying what happens if.
\end{grammarlens}

\subsection*{Your turn}
\begin{itemize}
\raggedright
  \item \emph{Say it:} \emph{paḷvāṭ}
  \item \emph{Say it:} \emph{paḷvāṭ}, once more
  \item \emph{Say it:} \emph{yā niymāt ek paḷvāṭ āhe.} — There is a loophole in this rule.
  \item \emph{Recall:} say the Marathi for to block, then the Marathi for to take notice, then say \emph{paḷvāṭ} again
\end{itemize}

\begin{tcolorbox}[breakable,colback=teal!4,colframe=teal!35!black,title={Before you move on}]
What does \mr{पळवाट} mean? (\textquotedblleft{}A loophole; an escape route\textquotedblright{}.) Say it once more.
\end{tcolorbox}
